% =====================================================================
% Hardware Implementation  --  TARGET: ~600 words
% Marking weight: 20 marks (the biggest single rubric row)
% Brief asks for: assembly, sensors + interfacing, schematics, power
% management & protection, hardware challenges.
% =====================================================================
\section{Hardware Implementation}
\label{sec:hardware}

\guide{
Twenty marks --- this section needs to be rich. The 1st+ rubric wants
``industrial-quality hardware design with strong reliability,
protection, and performance optimisation.'' Five sub-sections:
\begin{enumerate}
  \item Bill of materials \& physical assembly.
  \item Sensor interfacing (joystick axes, joint potentiometers, IR
        proximity for autonomous mode).
  \item Schematics (already drawn --- include figures for arm
        wiring, controller wiring, HC-05 divider).
  \item Power architecture details + decoupling, current-rail
        separation.
  \item Hardware challenges and how they were resolved.
\end{enumerate}
The last sub-section is where you score the ``robustness improvements''
language in the 1st+ rubric --- e.g.\ servo brown-out fixed by
adding bulk caps, ADC noise cleaned by RC filtering, etc.
}

\subsection{Bill of materials and physical assembly}
\placeholder{[\ldots one short paragraph listing the principal
components: Mega 2560, Uno R4 Minima, 3 $\times$ MG996R, 1 $\times$ SG90,
HC-05 pair, 4 $\times$ 10\,k$\Omega$ potentiometers, joysticks, IR
sensor. Then a paragraph on the mechanical assembly --- which member
built it, what tooling was used, what torque values for the servo
horns, and how cables were routed.\ldots]}

% TODO PHOTO: assembled arm on the bench with cabling tidy.
\begin{figure}[t]
  \centering
  \fbox{\parbox{0.92\columnwidth}{\centering\vspace{4pt}
    \emph{[Photograph of the assembled arm on the bench. Add once
    final assembly is complete.]}\vspace{4pt}}}
  \caption{Assembled arm prototype (photo placeholder).}
  \label{fig:armphoto}
\end{figure}

\subsection{Sensor interfacing}
\guide{For each sensor, write: working principle (1 sentence),
electrical interface to the MCU (pin / peripheral), and what was done
to clean the signal (filter, averaging, dead-band). This last bit is
where the rubric's ``correctly interfaced hardware'' language lives.}

\paragraph{Joystick axes}
\placeholder{[\ldots KY-023 thumbstick, two analog axes 0--3.3\,V,
ADC sampling at $\geq$200\,Hz, dead-band of $\pm$X counts to suppress
quiescent jitter\ldots]}

\paragraph{Joint potentiometers}
\placeholder{[\ldots 10\,k$\Omega$ linear pot per joint, geared or
directly coupled, RC low-pass on wiper, ADC sampled at 100\,Hz under
timer ISR, 16-sample boxcar average\ldots]}

\paragraph{IR proximity (autonomous mode)}
\placeholder{[\ldots short-range IR module gating the DETECT state of
the autonomous FSM\ldots]}

\subsection{Wiring schematics}
\guide{Refer to the three SVG-derived figures already in
\texttt{figures/}: arm wiring, controller wiring, HC-05 voltage
divider. One paragraph of prose around them, calling out the
voltage-divider on HC-05 RX as a deliberate design decision.}

\begin{figure}[t]
  \centering
  \includegraphics[width=0.95\columnwidth]{03_arm_wiring.pdf}
  \caption{Arm-side wiring: Mega 2560, four servos, HC-05 slave on
  Serial1, joint potentiometers on the analog inputs.}
  \label{fig:armwiring}
\end{figure}

\begin{figure}[t]
  \centering
  \includegraphics[width=0.95\columnwidth]{04_controller_wiring.pdf}
  \caption{Controller-side wiring: Uno R4 Minima, two KY-023
  joysticks, mode button, HC-05 master on Serial1.}
  \label{fig:ctrlwiring}
\end{figure}

\begin{figure}[t]
  \centering
  \includegraphics[width=0.65\columnwidth]{02_hc05_voltage_divider.pdf}
  \caption{HC-05 RX voltage divider: 2\,k$\Omega$/1\,k$\Omega$ drops
  the 5\,V Arduino TX to $\sim$3.3\,V at the HC-05 RX pin.}
  \label{fig:hc05div}
\end{figure}

\placeholder{[\ldots one paragraph of prose tying the three figures
together and noting the protection rationale\ldots]}

\subsection{Power and protection}
\placeholder{[\ldots one paragraph: split rails, decoupling caps,
common ground, fuse / polyswitch (if used), e-stop button wired into
the controller as an interrupt-driven flag\ldots]}

\subsection{Hardware challenges and resolutions}
\guide{This is the rubric's ``robustness improvements'' slot.
Concrete examples to cover:
\begin{itemize}
  \item Servo brown-out / Mega resetting under simultaneous motion ---
        fix: add 470\,$\mu$F bulk capacitor on the servo rail.
  \item ADC noise on potentiometer reads --- fix: RC filter +
        averaging.
  \item HC-05 RX 5\,V tolerance --- fix: voltage divider.
  \item Crosstalk on the breadboard --- fix: dedicated ground rail.
\end{itemize}
Each item is one or two sentences. Concrete numbers always score
better than generalities.}

\placeholder{[\ldots one paragraph or short bulleted list per
challenge\ldots]}

\corehit{Hardware brief Objectives 1, 3 and~5: hardware components
assembled, sensors interfaced, motor control wired up; circuit
schematics for both transmitter and receiver provided.}

\advhit{[\ldots e.g.\ analogue-front-end design (RC filter + digital
averaging on potentiometers); split-rail power architecture with
documented grounding strategy; voltage-level translation on HC-05
RX; hardware e-stop wired through firmware to a fail-safe pose.\ldots]}
